\chapter{KESIMPULAN DAN SARAN}

Bab ini menyajikan kesimpulan, saran, dan keterbatasan penelitian berdasarkan hasil analisis pada Bab IV. Kesimpulan disusun untuk menjawab rumusan masalah mengenai kondisi strategi pemasaran Nasi Gerilya pada GrabFood, faktor internal dan eksternal yang memengaruhi strategi, serta alternatif strategi pemasaran yang dapat digunakan untuk meningkatkan penjualan. Penarikan kesimpulan dikaitkan kembali dengan teori strategi pemasaran, bauran pemasaran jasa, pemasaran digital, dan SWOT/IFAS-EFAS agar rekomendasi yang diberikan tetap berhubungan dengan konsep akademik dan bukti lapangan \parencite{kotler2023principlesmarketing,wirtz2024servicesmarketing,chaffey2022digitalmarketing,david2023strategicmanagement}.

\section{Kesimpulan}
\begin{enumerate}
    \item Strategi pemasaran Rumah Makan Nasi Gerilya pada platform GrabFood sudah memiliki dasar yang kuat pada aspek produk, porsi, reputasi digital, penggunaan promo, dan kemauan berbenah melalui sistem digital. Menu yang paling menonjol dari data pelanggan, pegawai, dan intern \textit{business consultant} adalah dendeng, ayam pop, dan ayam goreng, dengan ayam pop layak diperkuat sebagai hero menu. Porsi nasi dinilai besar, kemasan relatif rapi, rating GrabFood berada pada posisi kuat, dan promo/voucher membantu menurunkan hambatan harga. Namun, strategi tersebut belum sepenuhnya stabil karena masih terdapat kelemahan pada proses layanan, terutama akurasi packing, item terpisah/berkuah yang rawan tertinggal, antrean lintas kanal saat ramai, update stok, serta persepsi harga yang mahal ketika tidak ada promo atau ongkir murah. Dengan demikian, strategi pemasaran perlu dipahami sebagai kombinasi nilai produk, bukti digital, dan kualitas proses layanan, bukan hanya promosi atau harga \parencite{kotler2023principlesmarketing,wirtz2024servicesmarketing,chaffey2022digitalmarketing}.

    \item Faktor internal utama Nasi Gerilya terdiri atas kekuatan rasa, menu unggulan, porsi besar, reputasi digital, adanya alur kerja checker-display-dapur, dan arah digitalisasi internal seperti sistem antrean, menu digital, stamp digital, dan absensi digital. Kelemahan internal yang paling penting adalah akurasi pesanan saat peak hour, kecepatan layanan, antrean dine in/ojek online/take away, update stok, ketergantungan pelanggan pada promo, kurang jelasnya beberapa informasi menu seperti isi paket, kuah, sambal, label, dan batas catatan khusus, serta belum optimalnya pengelolaan data pelanggan langsung. Dari sisi eksternal, peluang terbesar berasal dari perilaku pelanggan GrabFood yang membutuhkan makanan berat praktis, promo/voucher yang mendorong pembelian, ruang perbaikan tampilan/menu digital, dan program loyalitas. Ancaman utama berasal dari kompetitor dengan reputasi kuat, pembanding harga dan promo, biaya tambahan/ongkir, biaya platform, kenaikan bahan baku, sensitivitas harga, serta risiko ulasan negatif akibat human error. Pembacaan ini sesuai dengan kerangka SWOT yang memisahkan faktor yang dapat dikendalikan usaha dan faktor lingkungan yang harus direspons melalui strategi \parencite{david2023strategicmanagement,gurel2017swot,rangkuti2016analisisSWOT}.

    \item Hasil IFAS sebesar 2,50 dan EFAS sebesar 2,51 menunjukkan bahwa posisi Nasi Gerilya berada pada kondisi sedang: usaha memiliki kekuatan dan peluang yang cukup besar, tetapi perlu memperbaiki proses agar peluang digital tidak berubah menjadi keluhan pelanggan. Bobot IFAS/EFAS disusun dari kekuatan bukti, frekuensi kemunculan tema, dan dampak faktor terhadap penjualan atau pembelian ulang, sehingga angka tersebut berfungsi sebagai alat bantu prioritas strategi, bukan ukuran statistik. Berdasarkan hasil triangulasi, strategi prioritas yang paling layak diterapkan adalah memperkuat sistem antrean, checklist packing, dan SOP peak hour; membuat paket menu unggulan berbasis segmen pelanggan; memperbaiki tampilan digital dan menu digital; menyinkronkan promo dengan stok, kapasitas dapur, HPP, dan margin; serta mengaktifkan stamp digital/database pelanggan untuk mendorong repeat order. Strategi tersebut dibaca sebagai upaya yang berpotensi mendukung peningkatan penjualan melalui perbaikan keputusan pembelian, repeat order, dan pengalaman layanan, bukan sebagai bukti kuantitatif bahwa penjualan sudah meningkat \parencite{rangkuti2016analisisSWOT,david2023strategicmanagement,hendrizon2024madayacoffee}.
\end{enumerate}

\section{Saran}

Saran disusun sebagai penerjemahan temuan SWOT ke tindakan pemasaran yang dapat dikendalikan usaha. Karena jasa kuliner pada platform digital sangat ditentukan oleh proses, orang, bukti fisik/digital, dan pengelolaan nilai pelanggan, saran diarahkan pada perbaikan layanan sebelum perluasan promosi \parencite{wirtz2024servicesmarketing,kotler2023principlesmarketing,chaffey2022digitalmarketing}.

\subsection{Saran bagi Rumah Makan Nasi Gerilya}
\begin{enumerate}
    \item Nasi Gerilya perlu menerapkan sistem antrean lintas kanal dan checklist packing khusus GrabFood. Sistem antrean perlu membedakan alur dine in, ojek online, dan take away, sedangkan checklist packing perlu mencakup nomor pesanan, lauk utama, item terpisah, kuah, sambal, add-on, catatan khusus, alat makan, segel, dan label pesanan. Hal ini penting karena temuan utama menunjukkan bahwa antrean menumpuk, item kurang lengkap, dan lauk terpisah merupakan risiko layanan yang paling langsung memengaruhi kepuasan pelanggan. Dari sisi teori pemasaran jasa, perbaikan proses menjadi prasyarat agar kualitas layanan lebih konsisten \parencite{wirtz2024servicesmarketing}.

    \item Nasi Gerilya perlu memperkuat SOP peak hour, terutama pada jam makan siang, malam, akhir pekan, hari Minggu, hari raya, dan saat promo aktif. SOP tersebut sebaiknya mencakup pembagian tugas checker, display, dan dapur; update stok; estimasi tunggu driver; penggunaan sistem antrean; serta komunikasi cepat apabila ada menu yang belum ready.

    \item Nasi Gerilya dapat membuat paket menu unggulan berdasarkan segmen pelanggan, misalnya paket dendeng untuk pelanggan yang mencari rasa khas, paket ayam pop untuk pekerja kantor, paket ayam goreng untuk mahasiswa atau pelanggan yang sensitif harga, serta bundling minum/sambal/kuah. Paket ini perlu dikaitkan dengan promo agar harga akhir terasa lebih wajar. Strategi paket berbasis segmen membantu usaha menyelaraskan targeting, positioning, dan nilai yang dirasakan pelanggan \parencite{kotler2023principlesmarketing}.

    \item Tampilan toko GrabFood dan menu digital perlu diperbaiki dengan menonjolkan menu hero, memperjelas isi paket, informasi kuah dan sambal, status stok, foto menu yang konsisten, serta label atau segel pada kemasan. Perbaikan ini dapat mengurangi keraguan pelanggan sebelum checkout, memperjelas ekspektasi terhadap makanan yang diterima, dan membantu pelanggan melakukan pemesanan secara lebih mandiri.

    \item Promo sebaiknya tidak dijalankan hanya untuk menaikkan jumlah order, tetapi disiapkan bersama stok, kapasitas dapur, HPP, dan margin. Menu yang sedang promo perlu diinformasikan lebih awal kepada dapur agar tidak menimbulkan keterlambatan, menu habis, penurunan kualitas layanan, atau tekanan margin yang terlalu besar.

    \item Nasi Gerilya dapat mengaktifkan stamp digital dan database pelanggan untuk program loyalitas dine in dan take away. Program ini dapat membantu mengingatkan pelanggan lama, memberi insentif pembelian ulang, dan mengurangi ketergantungan penuh pada voucher platform. Penggunaan data pelanggan juga mendukung pemasaran digital yang lebih terukur karena repeat order dan respons promosi dapat dievaluasi dari waktu ke waktu \parencite{chaffey2022digitalmarketing,kotler2021marketingmanagement}.

    \item Pengelola perlu memantau data sederhana secara rutin, seperti jumlah pesanan, omzet, waktu pelayanan, jumlah antrean tertangani, jumlah item kurang, menu yang sering habis, waktu tunggu driver, komplain, rating/ulasan, repeat order, paket yang paling laku, efektivitas sistem, dan perubahan harga kompetitor. Pemantauan ini dapat menjadi dasar evaluasi mingguan untuk menjaga repeat order dan mencegah penurunan penjualan.
\end{enumerate}

\subsection{Saran bagi Peneliti Selanjutnya}
\begin{enumerate}
    \item Penelitian selanjutnya dapat menambah jumlah responden pelanggan agar variasi pengalaman pembelian lebih luas, terutama pelanggan yang pernah berhenti membeli atau pelanggan yang membandingkan langsung Nasi Gerilya dengan kompetitor tertentu.

    \item Penelitian berikutnya dapat memasukkan data pemilik atau pengelola utama dalam wawancara penelitian utama agar pembahasan mengenai biaya promo, margin, kebijakan harga, dan arah pengembangan usaha menjadi lebih kuat.

    \item Penelitian selanjutnya dapat membandingkan performa Nasi Gerilya pada GrabFood, GoFood, dan ShopeeFood agar strategi kanal digital tidak hanya dibaca dari satu platform.

    \item Penelitian lanjutan dapat menambah wawancara kompetitor yang belum terwakili, terutama kompetitor yang sudah muncul dalam observasi GrabFood seperti Pondok Gurih, Garuda, Paripurna, dan Dendeng Batokok, agar pembacaan ancaman kompetitif tidak hanya bertumpu pada observasi visual dan dua wawancara pembanding.

    \item Penelitian lanjutan dapat menggunakan data kuantitatif yang lebih rinci, seperti jumlah order harian, conversion rate, nilai promo, biaya komisi, HPP, waktu layanan, antrean tertangani, menu terlaris, data stamp digital, dan repeat order, untuk menguji dampak strategi yang direkomendasikan.
\end{enumerate}

\section{Keterbatasan Penelitian}
Penelitian ini memiliki beberapa keterbatasan. Pertama, penelitian berfokus pada strategi pemasaran Nasi Gerilya pada GrabFood, sehingga hasilnya tidak dimaksudkan untuk menggambarkan seluruh kanal penjualan digital maupun offline secara menyeluruh. Kedua, pembahasan peningkatan penjualan ditulis sebagai potensi strategis berdasarkan triangulasi kualitatif, bukan sebagai pengujian kausal bahwa rekomendasi strategi sudah menghasilkan kenaikan penjualan. Ketiga, jumlah responden pelanggan terbatas, sehingga temuan pelanggan digunakan sebagai indikasi pengalaman pembelian dan bahan triangulasi, bukan sebagai generalisasi statistik. Keempat, pembacaan kompetitor menggabungkan observasi platform dan wawancara pembanding, sehingga hasilnya menunjukkan pola persaingan utama, bukan pemetaan menyeluruh terhadap seluruh restoran Padang di Medan. Kelima, tampilan, promo, harga, rating, dan ulasan pada GrabFood dapat berubah dari waktu ke waktu, sehingga hasil observasi menggambarkan kondisi pada tanggal pengamatan.
